\section{Conexión con los conceptos del curso}
\label{sec:curso}

El artículo \cite{piperno2025} se deja leer como una cadena de decisiones que el curso trata por separado. Cada etapa del flujo de procesamiento activa un concepto distinto y la correspondencia ayuda a ver por qué el método funciona y dónde pierde sustento.

Las representaciones contextuales son el sustrato. BERT produce un vector distinto para cada aparición de una palabra, y mBERT comparte un vocabulario de subpalabras entre lenguas, lo que permite que el mismo encoder procese español, francés y alemán \cite{devlin2019}. BioBERT conserva la arquitectura y desplaza la geometría hacia el vocabulario biomédico \cite{lee2020}. Que profesor y estudiante pertenezcan a la misma familia es una condición del diseño, porque una pérdida de distancia entre embeddings exige espacios comparables. El artículo usa [CLS] en el RSA, una decisión discutible: [CLS] se entrena para la predicción de la oración siguiente, no para similitud, y Sentence-BERT suele usar la media \cite{reimers2019}. El texto no especifica el pooling de la destilación.

La transferencia y la adaptación de dominio explican la jerarquía de profesores. mBERT representa el preentrenamiento general, BioBERT la adaptación al dominio y los modelos ajustados la adaptación a la tarea. La curva del 10 al 100\,\% es un experimento de eficiencia de datos. Una alternativa que el artículo no prueba es continuar el preentrenamiento de mBERT con texto biomédico en español antes de destilar: atacaría el problema desde el estudiante y permitiría comparar qué rinde más.

La destilación de características da nombre al método. El estudiante imita representaciones, no logits; la cabeza de clasificación del profesor no se transfiere. Eso explica por qué el estudiante puede ajustarse después con etiquetas en la lengua meta. También explica el alcance de la evidencia: el RSA verifica que la geometría se acercó, no que el contenido clínico se conservó. Una alternativa sería destilar sobre los 18 logits del profesor o combinar ambas señales.

El aprendizaje contrastivo aclara la pérdida. MNRL es la entropía cruzada de una softmax sobre similitudes: el positivo es la contraparte y los negativos son los demás pares del lote. El tamaño de lote define la dificultad y $\alpha$ cumple un papel parecido a la temperatura. MSE obliga a igualar coordenadas; la pérdida contrastiva solo exige que el positivo quede por delante. Una alternativa sería el minado de negativos difíciles o una pérdida contrastiva supervisada con las 18 clases.

La clasificación y el F1 cierran el ciclo. El ajuste final es una clasificación de 18 clases, y el F1 es el resumen. Macro, micro y ponderado responden preguntas distintas: macro promedia las clases por igual, micro acumula aciertos y ponderado pesa por soporte. Sin esa definición, el número de la Tabla~\ref{tab:f1} no se puede interpretar. Las matrices de confusión ofrecen la vista por clase que el promedio oculta, y el artículo las usa para comparar estudiantes, no para analizar errores. Una alternativa sería reportar F1 macro con soporte por clase y un bootstrap pareado.

Un enfoque lingüístico o de reglas sobre léxico terminológico, con la CIE-10 o el CUPS, sería más interpretable y auditable, pero más frágil ante variación, abreviaturas y texto libre. El artículo no lo compara; su ausencia confirma que la comparación es entre variantes de aprendizaje profundo, no entre familias de métodos.
